NLP benchmarks are routinely ``solved''---performance saturates near human parity---yet no automated system detects when this saturation occurs, leaving benchmark retirement decisions to informal community consensus that can lag the actual event by months.
We propose an automated saturation detection pipeline that reframes this problem as logistic curve fitting: given a benchmark's score-over-time leaderboard timeseries, we fit a bounded logistic growth model and formally compare it against linear and power-law alternatives via the Akaike Information Criterion.
Applied to GLUE and SuperGLUE, the logistic model is overwhelmingly preferred ($\Delta\text{AIC} \approx -194$ to $-250$ versus linear and $-113$ to $-357$ versus power-law alternatives---decisive by any standard threshold), and a dual-criterion saturation detector recovers the community-recognized saturation events within 3 and 5 months respectively.
A surprising finding emerges: the logistic inflection point predates benchmark launch for both benchmarks, suggesting that BERT-era pretraining had already saturated benchmark-exploitable signal before public leaderboard tracking began.
We further show that the pipeline generalizes reliably to benchmarks with as few as 30 leaderboard entries with bounded fitting.
Together, these results establish, to our knowledge, the first quantitative, reproducible foundation for automated benchmark saturation monitoring---transforming what has been a community judgment call into a principled, data-driven measurement.
